\scanpage{420}
\mspage{159}
\unsure{semblable|similaire} déjà connu pour $\alpha_2$, ces trois foncteurs sont des Hot-équivalences. Donc il en est de $\unsure{=|\text{même}}$ \unsure{de|des} $\unsure{\pi|\Pi}_1,\pi_2$. Comme $p_1,p_0,p_2$ sont des Hot-fibrations (étant propres et lisses), il s'ensuit que $\pi_1$ et $\pi_2$ induisent des Hot équivalences sur les fibres
\diagram{\begin{tikzcd}[ampersand replacement=\&,column sep=large]
\underline{\mathrm{Ch}}(X;a,b) \& \mathcal C_{\mathbb N^{\unsure{o|0}}}(a,b)\arrow[l,"{\pi_1(a,b)}"']\arrow[r,"{\pi_2(a,b)}"] \& \underline{\mathrm{Ch}}_\infty(X;a,b).
\end{tikzcd}}
Je l'avais vérifié (sans mal) déjà pour $\pi_2$\inserted{ et les $\pi_2(a,b)$}, \unsure{ce} \deleted{\unsure{la chose de}}\inserted{qui avait} permis, en inversant $\pi_2$ dans Hot, de définir \uline{dans Hot}
\[
\left\{\begin{aligned}
\underline{\mathrm{Ch}}_\infty(X)&\mathrel{\unsure{\text{---}|\longrightarrow}}\underline{\mathrm{Ch}}(X)\\
\underline{\mathrm{Ch}}_\infty(X;a,b)&\longrightarrow\underline{\mathrm{Ch}}(X;a,b),
\end{aligned}\right.
\]
et dire que ce sont des iso signifie justement que $\pi_1$ et les $\pi_1(a,b)$ sont des iso.

Je ne crois pas que $\mathcal C_{\mathbb N^o}$ ait en lui-même d'autre intérêt, que de permettre de comparer les $X\times X$-catégories $\underline{\mathrm{Ch}}_\infty(X)$ et $\underline{\mathrm{Ch}}(X)$, à l'aide du diagramme précédent. Par contre, $\underline{\mathrm{Ch}}_\infty(X)$ et $\underline{\mathrm{Ch}}(X)$ m'apparaissent comme les deux ingrédients essentiels du formalisme des chemins dans \unsure{(Cat)|Cat}.
% archive sheet number: 419
